\clearpage
\chapter{SYSTEM DESIGN AND GLOBAL ARCHITECTURE}

\section{Introduction}
This chapter translates the requirements and selected technologies into a coherent system design. It describes data contracts, layer boundaries, component interactions, deployment roles, workflow dependencies, and security boundaries without repeating implementation commands or evaluation results.

\section{Conceptual Data Design}

\subsection{Active Data Sources}
Client~1 contains three complementary sources:
\begin{itemize}
    \item \textbf{Bank Marketing}, representing customer attributes, campaign contacts, outcomes, and macro-economic context;
    \item \textbf{Online Shoppers Purchasing Intention}, representing web-session engagement and purchase-intention signals;
    \item \textbf{Online Retail II}, representing invoice, product, revenue, and repeat-purchase behavior.
\end{itemize}
Together they provide campaign, behavioral-intent, and commercial dimensions suitable for a Customer~360 analytical foundation.

\subsection{Customer 360 Modeling Principle}
The platform does not claim deterministic enterprise identity resolution between unrelated academic datasets. Instead, each source is standardized within its own semantic domain and contributes features to a customer-oriented analytical foundation. The design supports conversion, engagement, activity, and value analysis while preserving the boundary between implemented engineering and future ML products.

\subsection{Lakehouse Layering Strategy}
\begin{table}[htbp]
    \centering
    \small
    \arrayrulecolor{ModernTableRule}
    \begin{tabularx}{\textwidth}{|p{0.19\textwidth}|p{0.29\textwidth}|X|}
        \hline
        \rowcolor{ModernTableHeader}
        \color{white}\textbf{Layer} & \color{white}\textbf{Responsibility} & \color{white}\textbf{Design rule} \\
        \hline
        Bronze & Replayable source landing & Preserve source identity and separate ingestion from processing \\
        \hline
        Raw Iceberg & Parsed governed source tables & Apply explicit schemas while retaining source-level meaning \\
        \hline
        Staging & Standardization & Rename, type, clean, and expose consistent source fields \\
        \hline
        Intermediate & Reusable domain logic & Aggregate contacts, sessions, sales, and feature-oriented signals \\
        \hline
        Analytics & Consumption-ready products & Present Customer~360 overview, conversion, and sales outputs \\
        \hline
    \end{tabularx}
    \caption{Conceptual lakehouse layers}
    \label{tab:design-lakehouse-layers}
    \arrayrulecolor{black}
\end{table}

\section{Global Platform Architecture}

\subsection{End-to-End Data Flow}
The authoritative flow is:
\begin{center}
\begin{tikzpicture}[node distance=7mm and 7mm]
    \node[workflow/step,fill=WorkflowBlue!15,text width=2.6cm] (source) {Customer\\datasets};
    \node[workflow/step,fill=WorkflowOrange!17,text width=2.2cm,right=of source] (kafka) {Kafka\\ingestion};
    \node[workflow/step,fill=WorkflowTeal!16,text width=2.2cm,right=of kafka] (hdfs) {HDFS\\bronze};
    \node[workflow/step,fill=WorkflowPurple!14,text width=2.2cm,right=of hdfs] (spark) {Spark\\raw load};
    \draw[workflow/arrow] (source) -- (kafka);
    \draw[workflow/arrow] (kafka) -- (hdfs);
    \draw[workflow/arrow] (hdfs) -- (spark);

    \node[workflow/step,fill=WorkflowBlue!13,text width=2.3cm,below=12mm of spark] (iceberg) {Iceberg +\\Metastore};
    \node[workflow/step,fill=WorkflowTeal!14,text width=2.2cm,left=of iceberg] (trino) {Trino\\SQL};
    \node[workflow/step,fill=WorkflowOrange!14,text width=2.2cm,left=of trino] (dbt) {dbt-Spark\\models};
    \node[workflow/step,fill=WorkflowPurple!12,text width=2.3cm,left=of dbt] (quality) {GX quality\\evidence};
    \draw[workflow/arrow] (spark) -- (iceberg);
    \draw[workflow/arrow] (iceberg) -- (trino);
    \draw[workflow/arrow] (trino) -- (dbt);
    \draw[workflow/arrow] (dbt) -- (quality);
\end{tikzpicture}
\end{center}
Airflow coordinates the operational order, while Prometheus and Grafana observe infrastructure and pipeline state. The diagram expresses data responsibility rather than physical network placement.

\subsection{Component Interaction}
Kafka decouples source readers from bronze writers. HDFS preserves landed files independently from source systems. Spark converts bronze objects into Iceberg tables registered through Hive Metastore. Spark and Trino share the catalog, while dbt connects through Spark Thrift to construct curated models. GX validates raw lakehouse data and retains evidence. Airflow invokes the stages, and the monitoring stack collects the resulting operational signals.

\subsection{Metadata and Catalog Design}
Hive Metastore acts as the metadata authority for raw, staging, intermediate, and analytics namespaces. Iceberg stores table-level metadata and snapshots. This combination separates physical HDFS files from governed table state and allows Spark and Trino to interpret the same objects consistently.

\subsection{Transformation and Quality Design}
The dbt graph contains three staging models, four intermediate models, and three analytics models. GX contains one raw-lakehouse validation definition for each active dataset. Quality is treated as a gate with retained artifacts rather than a visual dashboard status alone.

\section{Distributed Deployment Design}

\subsection{Local Control Plane}
The local workstation contains orchestration, monitoring, user interfaces, dbt and GX runtimes, and the client used to initiate remote Spark work. Keeping these services local protects the limited VM resources and provides convenient access during operation and demonstration.

\subsection{Three-VM Data Plane}
VM2 is the leader node because it hosts the HDFS namespace authority, shared metastore, Spark master and SQL services, and Trino coordinator. VM1 and VM3 extend ingestion, storage, processing, and querying through broker, DataNode, Spark-worker, and Trino-worker services.

\subsection{Logical Service Placement}
\begin{table}[htbp]
    \centering
    \footnotesize
    \arrayrulecolor{ModernTableRule}
    \begin{tabularx}{\textwidth}{|>{\RaggedRight\arraybackslash}p{0.12\textwidth}|>{\RaggedRight\arraybackslash}p{0.17\textwidth}|>{\RaggedRight\arraybackslash}X|}
        \hline
        \rowcolor{ModernTableHeader}
        \color{white}\textbf{Node} & \color{white}\textbf{Design role} & \color{white}\textbf{Service responsibilities} \\
        \hline
        Local PC & Control plane & Airflow, Prometheus, Grafana, Kafka UI, pipeline and PostgreSQL exporters, dbt, GX, and remote Spark orchestration \\
        \hline
        VM1 & Worker & Kafka broker 1, HDFS DataNode 1, Spark worker 1, Trino worker 1, and cAdvisor \\
        \hline
        VM2 & Leader and worker & Kafka broker 2; HDFS NameNode and DataNode 2; metastore service and database; Spark master, worker 2, history, Thrift, and submit services; Trino coordinator; cAdvisor \\
        \hline
        VM3 & Worker & Kafka broker 3, HDFS DataNode 3, Spark worker 3, Trino worker 2, and cAdvisor \\
        \hline
    \end{tabularx}
    \caption{Logical control-plane and data-plane placement}
    \label{tab:design-logical-placement}
    \arrayrulecolor{black}
\end{table}

\subsection{Remote Spark Execution Model}
Airflow reaches VM2 through SSH and invokes the Spark submit helper located with the cluster. The Spark master then allocates work across the available workers. This avoids executor-to-driver communication problems caused by a dynamically addressed local driver and keeps raw processing within the data plane.

\section{Workflow Design}

\subsection{Infrastructure Startup Sequence}
The designed startup order is VM2, VM1, VM3, and the control plane. Leader services become available first, workers then register with them, and orchestration starts only after the data plane is reachable.

\subsection{Five-DAG Operational Sequence}
\begin{enumerate}
    \item environment reset;
    \item lakehouse setup;
    \item lakehouse readiness;
    \item Kafka--HDFS--Spark lakehouse construction and GX validation;
    \item dbt-Spark transformation and testing.
\end{enumerate}
This order separates destructive cleanup, initialization, readiness, raw processing, and curated transformation into visible failure boundaries.

\subsection{Failure Handling and Reprocessing}
A failed stage remains recorded in Airflow. Reset utilities target controlled runtime state, while stage separation permits rerunning a failed boundary without presenting historical failure as success. Source, table, and model naming contracts remain stable across reruns.

\section{Security and Operational Design}
The design uses explicit node separation, known service endpoints, environment-scoped secrets, controlled SSH submission, and centralized user-facing interfaces. Private keys and secret values are operational inputs and must never appear in report screenshots, code excerpts, or appendices. The design is a baseline secured deployment posture, not a zero-trust or enterprise IAM architecture.

\section{Conclusion}
This chapter defined how the selected stack satisfies the requirements through clear data, workflow, and deployment boundaries. Chapter~VI documents the concrete repository, node specifications, services, scripts, and operational execution used to realize this design.